\begin{abstract}
In the completed observed-catalogue tests, the tested stress-depletion additions do not improve retrospective earthquake prediction: they lose to matched excitation models or select zero inhibition. Standard unweighted inverse-power ETAS outperforms the target background-weighted offspring model in every reported California and Japan primary test fold. These negative procedure-level results motivate a restricted mathematical and empirical assessment of inhibition, rather than a claim of physical falsification.

The analysis distinguishes autonomous physical loading from diagnostic laws with independently specified backgrounds, finite exponential mixtures from exact fractional memory, and full-law identification from finite-catalogue estimation. Under explicit hypotheses, the finite-mode process has first-moment stability while stationary second moments may diverge. Exact-memory obstructions concern specified no-feedback submodels. Stress-memory parameters disappear at zero inhibition. Finite-history likelihood tools do not supply general boundary asymptotics.

Controlled studies show sensitivity to estimated background nuisance structure, but generating-parameter null calibration is an oracle design, and the later four-procedure continuation now has completed missing-only recovery with full planned-denominator frequencies. This remains an oracle calibration, not a fitted-null plug-in study or a universal size-matched power result. New Zealand model development repeatedly reused the same held-out events; freezing each later protocol did not restore untouched-test validity. Its successive comparisons are exploratory. Large continuous-location gains mainly repair coarse uniform cell-lift references, and do not establish superiority over competitive spatial models. The original California timing gains concentrated in one retained sequence; the expanded causal-window backtest gives variable procedure contrasts; Japan reverses the matched timing direction and fails the registered combined transfer criterion. Fixing the Japanese background trend at zero reduces but does not remove the timing reversal. Conditional simulations reveal count, temporal and spatial discrepancies. The completed continuous-statistic within-origin omnibus checks flag both OC and EP; the retained minimum-rank procedure is less discriminating in these data.

Continuous immigrant-only inhibition, including the magnitude-scaled LI+ revision, selects zero strength in every observed California fold. Raw total-intensity LI+ loses held-out joint score in every fold, while its complexity-gated version selects the uninhibited floor. Its historical eight fixed-seed controls included four null losses and two narrow/short positive-inhibition losses, alongside two broad/long gains. The subsequent expanded fixed-backbone matrix evaluates magnitude-scaled immigrant-only and total-intensity inhibition. Magnitude-scaled immigrant-only inhibition gains in the broad/long synthetic regime while total-intensity attenuation has a negative mean gain there. The matrix also retains negative outcomes and shows that penalized selection can remove null overfitting while missing narrow/short positive inhibition. Neither study establishes general detection or physical causality. Unit stress marks and unattenuated triggering constrain the original LI model. The conclusions concern declared fitted procedures on inspected historical data, not the existence or absence of physical inhibition. The prospective prototype remains disabled and supplies no forecast score.
\end{abstract}
